\documentclass{article}
\usepackage{amsmath,amssymb,amsthm}
\usepackage{algorithm}
\usepackage{algpseudocode}
\usepackage{enumitem}
\usepackage{hyperref}
\usepackage{bm}

\newtheorem{theorem}{Theorem}
\newtheorem{lemma}{Lemma}
\newtheorem{assumption}{Assumption}

\begin{document}

\title{Trust‑Region Method for Convex \& Smooth Optimization \\ Detailed Algorithmic Description and Convergence Proof}
\author{}
\date{}
\maketitle

%%%%%%%%%%%%%%%%%%%%%%%%%%%%%%%%%%%%%%%%%%%%%%%%%%%%%%%%%%%%%%%%%%%%%%%%%%%%%%%%
\section*{Algorithm Name}
\textbf{Convex Trust‑Region (CTR) Method}

%%%%%%%%%%%%%%%%%%%%%%%%%%%%%%%%%%%%%%%%%%%%%%%%%%%%%%%%%%%%%%%%%%%%%%%%%%%%%%%%
\section*{Mathematical Formulation}
Let $f:\mathbb{R}^{n}\rightarrow\mathbb{R}$ be a convex, twice continuously differentiable function.
At iteration $k$ we have the current point $x_{k}\in\mathbb{R}^{n}$ and a trust‑region radius
$\Delta_{k}>0$.

\medskip
\noindent\textbf{Quadratic model.}
\[
m_{k}(p)\;:=\;f(x_{k})+\nabla f(x_{k})^{\top}p+\tfrac12 p^{\top}B_{k}p,
\qquad B_{k}:=\nabla^{2}f(x_{k}) .
\]

\noindent\textbf{Subproblem (trust‑region).}
\[
\min_{p\in\mathbb{R}^{n}}\; m_{k}(p)\quad\text{s.t.}\quad \|p\|_{2}\le \Delta_{k}.
\tag{TR$_k$}
\]

We assume that (TR$_k$) is solved \emph{approximately} obtaining a step $p_{k}$ that satisfies
the \emph{Cauchy decrease condition}
\[
m_{k}(0)-m_{k}(p_{k})\;\ge\;\kappa_{\mathrm{c}}\,
\bigl(m_{k}(0)-m_{k}(p^{\mathrm{C}}_{k})\bigr),\qquad
\kappa_{\mathrm{c}}\in(0,1],
\tag{1}
\]
where $p^{\mathrm{C}}_{k}$ is the exact Cauchy point (steepest‑descent step projected onto the
trust region).

\medskip
\noindent\textbf{Actual vs. predicted reduction.}
\[
\begin{aligned}
\text{Actual reduction: }&\rho^{\mathrm{act}}_{k}
:= f(x_{k})-f(x_{k}+p_{k}),\\[2mm]
\text{Predicted reduction: }&\rho^{\mathrm{pred}}_{k}
:= m_{k}(0)-m_{k}(p_{k}),\\[2mm]
\text{Ratio: }&\rho_{k}:=
\frac{\rho^{\mathrm{act}}_{k}}{\rho^{\mathrm{pred}}_{k}} .
\end{aligned}
\tag{2}
\]

\noindent\textbf{Acceptance and radius update.}
Given constants $0<\eta_{1}<\eta_{2}<1$ and $0<\gamma_{\mathrm{inc}}>1$, $0<\gamma_{\mathrm{dec}}<1$:

\[
x_{k+1}= 
\begin{cases}
x_{k}+p_{k}, & \text{if }\rho_{k}\ge \eta_{1}\quad\text{(accept step)},\\[1mm]
x_{k},      & \text{otherwise (reject)} .
\end{cases}
\tag{3}
\]

\[
\Delta_{k+1}= 
\begin{cases}
\min\{\gamma_{\mathrm{inc}}\Delta_{k},\;\Delta_{\max}\}, & \text{if }\rho_{k}\ge \eta_{2},\\[1mm]
\gamma_{\mathrm{dec}}\Delta_{k}, & \text{if }\rho_{k}<\eta_{1},\\[1mm]
\Delta_{k}, & \text{otherwise}.
\end{cases}
\tag{4}
\]

The algorithm terminates when $\| \nabla f(x_{k})\| \le \varepsilon$ or after a maximum
number of iterations.

%%%%%%%%%%%%%%%%%%%%%%%%%%%%%%%%%%%%%%%%%%%%%%%%%%%%%%%%%%%%%%%%%%%%%%%%%%%%%%%%
\section*{Pseudocode}
\begin{algorithm}[H]
\caption{Convex Trust‑Region (CTR) Method}
\begin{algorithmic}[1]
\State \textbf{Input:} initial point $x_{0}$, initial radius $\Delta_{0}>0$,
                 parameters $\eta_{1},\eta_{2},\gamma_{\mathrm{inc}},\gamma_{\mathrm{dec}}$, 
                 tolerance $\varepsilon$, max‑iter $K_{\max}$.
\For{$k=0,1,\dots,K_{\max}$}
    \State Compute $g_{k}\gets \nabla f(x_{k})$, $B_{k}\gets\nabla^{2}f(x_{k})$.
    \State Approximately solve (TR$_k$) to obtain $p_{k}$ satisfying (1).
    \State Compute $\rho_{k}$ via (2).
    \If{$\rho_{k}\ge \eta_{1}$}
        \State $x_{k+1}\gets x_{k}+p_{k}$ \Comment{Accept}
    \Else
        \State $x_{k+1}\gets x_{k}$ \Comment{Reject}
    \EndIf
    \State Update $\Delta_{k+1}$ by (4).
    \If{$\|g_{k}\|\le \varepsilon$}
        \State \textbf{break}
    \EndIf
\EndFor
\State \textbf{Output:} $x_{\text{out}}:=x_{k}$.
\end{algorithmic}
\end{algorithm}

%%%%%%%%%%%%%%%%%%%%%%%%%%%%%%%%%%%%%%%%%%%%%%%%%%%%%%%%%%%%%%%%%%%%%%%%%%%%%%%%
\section*{Dry Run Example}
Consider the one‑dimensional quadratic
\[
f(w)=(w-3)^{2},\qquad \nabla f(w)=2(w-3),\qquad \nabla^{2}f(w)=2 .
\]
Choose
\[
w_{0}=0,\;\Delta_{0}=0.5,\;
\eta_{1}=0.25,\;\eta_{2}=0.75,\;
\gamma_{\mathrm{inc}}=2,\;\gamma_{\mathrm{dec}}=0.5,
\]
and use the exact solution of (TR$_k$) (which coincides with the Cauchy point for a quadratic).

\medskip
\noindent\textbf{Iteration 0}
\[
g_{0}=2(0-3)=-6,\quad B_{0}=2.
\]
Exact minimizer of the quadratic model subject to $|p|\le0.5$:
\[
p_{0}^{*}= -\frac{g_{0}}{B_{0}} = 3,\quad\text{projected: }p_{0}= \operatorname{sign}(-g_{0})\min\{3,\Delta_{0}\}=0.5 .
\]
Predicted reduction:
\[
\rho^{\mathrm{pred}}_{0}=m_{0}(0)-m_{0}(p_{0})
=0-\bigl(g_{0}p_{0}+\tfrac12 B_{0}p_{0}^{2}\bigr)
= -(-6)(0.5)-\tfrac12(2)(0.5)^{2}=3-0.25=2.75 .
\]
Actual reduction:
\[
f(w_{0})=9,\qquad f(w_{0}+p_{0})=f(0.5)=(0.5-3)^{2}=6.25,
\]
\[
\rho^{\mathrm{act}}_{0}=9-6.25=2.75.
\]
Thus $\rho_{0}=1\ge\eta_{1}$ (accept) and $\rho_{0}\ge\eta_{2}$ (increase radius):
\[
w_{1}=0.5,\qquad \Delta_{1}= \min\{2\cdot0.5,\Delta_{\max}\}=1.0 .
\]

\medskip
\noindent\textbf{Iteration 1}
\[
g_{1}=2(0.5-3)=-5,\; B_{1}=2.
\]
Unconstrained step $p^{*}_{1}= -g_{1}/B_{1}=2.5$, but $\Delta_{1}=1$, so $p_{1}=1$.
\[
\rho^{\mathrm{pred}}_{1}= -g_{1}p_{1}-\tfrac12 B_{1}p_{1}^{2}=5\cdot1-1=4.
\]
\[
f(w_{1})=6.25,\qquad f(w_{1}+p_{1})=f(1.5)=(1.5-3)^{2}=2.25,
\]
\[
\rho^{\mathrm{act}}_{1}=6.25-2.25=4,\qquad \rho_{1}=1\ge\eta_{2}.
\]
Accept step, increase radius:
\[
w_{2}=1.5,\qquad \Delta_{2}= \min\{2\cdot1,\,\Delta_{\max}\}=2 .
\]

\medskip
\noindent\textbf{Iteration 2}
\[
g_{2}=2(1.5-3)=-3,\; B_{2}=2,\quad p^{*}_{2}=1.5.
\]
Now $\Delta_{2}=2$, so the unconstrained step is feasible: $p_{2}=1.5$.
\[
\rho^{\mathrm{pred}}_{2}= -g_{2}p_{2}-\tfrac12 B_{2}p_{2}^{2}=3\cdot1.5- \tfrac12\cdot2\cdot(1.5)^{2}=4.5-2.25=2.25.
\]
\[
f(w_{2})=2.25,\qquad f(w_{2}+p_{2})=f(3)=0,
\]
\[
\rho^{\mathrm{act}}_{2}=2.25,\qquad \rho_{2}=1\ge\eta_{2}.
\]
Accept step; the algorithm has reached the global minimiser $w^{*}=3$ after three
iterations.

%%%%%%%%%%%%%%%%%%%%%%%%%%%%%%%%%%%%%%%%%%%%%%%%%%%%%%%%%%%%%%%%%%%%%%%%%%%%%%%%
\section*{Assumptions}
\begin{assumption}[Convexity and Smoothness]\label{as:convex}
$f:\mathbb{R}^{n}\to\mathbb{R}$ is convex and twice continuously differentiable.
Moreover, its gradient is $L$‑Lipschitz:
\[
\|\nabla f(x)-\nabla f(y)\|\le L\|x-y\|,\qquad\forall x,y\in\mathbb{R}^{n}.
\tag{A1}
\]
\end{assumption}

\begin{assumption}[Bounded Hessian]\label{as:hessian}
There exist constants $0<\mu\le M$ such that
\[
\mu I \preceq \nabla^{2}f(x) \preceq M I,\qquad\forall x\in\mathbb{R}^{n}.
\tag{A2}
\]
\end{assumption}

\begin{assumption}[Approximate Subproblem Solution]\label{as:subprob}
The step $p_{k}$ returned by the subproblem solver satisfies the Cauchy decrease
condition (1) with a fixed constant $\kappa_{\mathrm{c}}\in(0,1]$.
\tag{A3}
\end{assumption}

\begin{assumption}[Parameter Bounds]\label{as:params}
\[
0<\eta_{1}<\eta_{2}<1,\qquad
0<\gamma_{\mathrm{dec}}<1<\gamma_{\mathrm{inc}},\qquad
\Delta_{\max}>0.
\tag{A4}
\]
\end{assumption}

All subsequent results are derived under \ref{as:convex}–\ref{as:params}.

%%%%%%%%%%%%%%%%%%%%%%%%%%%%%%%%%%%%%%%%%%%%%%%%%%%%%%%%%%%%%%%%%%%%%%%%%%%%%%%%
\section*{Main Convergence Theorem}
\begin{theorem}[Global Convergence and Rate]\label{thm:main}
Let $\{x_{k}\}$ be the sequence generated by the Convex Trust‑Region method.
Assume \ref{as:convex}–\ref{as:params} and that the subproblem is solved
in the sense of \ref{as:subprob}.
Define $f^{*}:=\inf_{x}f(x)$. Then for any $K\ge1$,
\[
\min_{0\le k<K}\|\nabla f(x_{k})\|^{2}
\;\le\;
\frac{2L\bigl(f(x_{0})-f^{*}\bigr)}
{\kappa_{\mathrm{c}}\,(1-\eta_{1})\,\eta_{1}\,K}.
\tag{5}
\]
Consequently,
\[
\lim_{k\to\infty}\|\nabla f(x_{k})\|=0,
\qquad
\text{and}\qquad
\|\nabla f(x_{k})\|=O\!\bigl(k^{-1/2}\bigr).
\]
\end{theorem}

\begin{theorem}[Finite‑Termination under Exact Solution]\label{thm:exact}
If the subproblem is solved exactly (i.e., $p_{k}$ is the global minimiser of (TR$_k$)),
then after at most
\[
\Bigl\lceil \frac{2L\bigl(f(x_{0})-f^{*}\bigr)}{\eta_{1}(1-\eta_{2})\Delta_{0}^{2}} \Bigr\rceil
\]
iterations the algorithm terminates with $\| \nabla f(x_{k})\|\le\varepsilon$.
\end{theorem}

%%%%%%%%%%%%%%%%%%%%%%%%%%%%%%%%%%%%%%%%%%%%%%%%%%%%%%%%%%%%%%%%%%%%%%%%%%%%%%%%
\section*{Theoretical Properties}
\begin{enumerate}[label=\arabic*)]
\item \textbf{Stability.}  The trust‑region radius $\Delta_{k}$ is never increased beyond
$\Delta_{\max}$ and is reduced whenever the model is inaccurate
($\rho_{k}<\eta_{1}$).  This guarantees that the iterates remain in a region where the
quadratic model is a good approximation of $f$ (by Lipschitz gradient).

\item \textbf{Hyperparameter Sensitivity.}
The constants $\eta_{1},\eta_{2}$ control the acceptance/rejection thresholds.
Larger $\eta_{1}$ yields more conservative steps (potentially more reductions of $\Delta_{k}$)
but improves the bound in \eqref{5} through the factor $(1-\eta_{1})\eta_{1}$.
The expansion factor $\gamma_{\mathrm{inc}}$ affects how fast the radius can grow; any
choice $>1$ maintains the $O(1/k^{1/2})$ rate as long as \ref{as:params} holds.

\item \textbf{Comparison to Gradient Descent.}
For convex $L$‑smooth $f$, plain gradient descent with stepsize $1/L$ attains
$f(x_{k})-f^{*}\le \frac{L\|x_{0}-x^{*}\|^{2}}{2k}$.
The trust‑region method enjoys a comparable \emph{gradient‑norm} rate (5) without
requiring a globally fixed stepsize; instead the radius adapts automatically to the
local curvature.

\item \textbf{Complexity per Iteration.}
Each iteration requires evaluating $\nabla f(x_{k})$, $\nabla^{2}f(x_{k})$ (or a
positive‑definite approximation) and solving a (generally) small quadratic
subproblem.  When $B_{k}$ is diagonal or low‑rank, the subproblem can be solved in
$O(n)$ time.

\end{enumerate}

\section*{Discussion}
The proof of Theorem~\ref{thm:main} follows the classical trust‑region analysis
but is specialised to the convex and $L$‑smooth setting, which permits a
simpler bound on the model error via the Lipschitz constant $L$.  The
Cauchy‑decrease condition \eqref{1} guarantees that even an approximate solution
provides a sufficient reduction proportionate to the gradient norm, which is the
key ingredient for deriving the $O(k^{-1/2})$ gradient‑norm rate.

%%%%%%%%%%%%%%%%%%%%%%%%%%%%%%%%%%%%%%%%%%%%%%%%%%%%%%%%%%%%%%%%%%%%%%%%%%%%%%%%
\section*{Preliminaries (Definitions and Lemmas)}
\begin{lemma}[Model Error]\label{lem:model}
For any $p$ with $\|p\|\le\Delta$,
\[
\bigl|\,f(x+p)-m_{k}(p)\,\bigr|\le \tfrac12 L\|p\|^{2}.
\tag{L1}
\]
\end{lemma}
\begin{proof}
By the $L$‑Lipschitz gradient property,
\[
f(x+p)=f(x)+\nabla f(x)^{\top}p+\int_{0}^{1}(1-t)\,p^{\top}\nabla^{2}f(x+tp)p\,dt .
\]
Using $ \|\nabla^{2}f(\cdot)\|\le L$ (from \ref{as:convex}) gives
\[
\bigl|f(x+p)-\bigl(f(x)+\nabla f(x)^{\top}p\bigr)\bigr|
\le \tfrac12 L\|p\|^{2}.
\]
Since $m_{k}(p)$ adds the term $\tfrac12 p^{\top}B_{k}p$ with $B_{k}=\nabla^{2}f(x)$,
the same bound holds for $|f(x+p)-m_{k}(p)|$.
\end{proof}

\begin{lemma}[Cauchy Decrease]\label{lem:cauchy}
Let $p^{\mathrm{C}}_{k}= -\frac{\Delta_{k}}{\|g_{k}\|}g_{k}$ be the Cauchy point.
Then
\[
m_{k}(0)-m_{k}(p^{\mathrm{C}}_{k})\;\ge\;
\frac{\Delta_{k}}{2}\,\|g_{k}\| -\frac{M}{2}\Delta_{k}^{2},
\tag{L2}
\]
where $M$ is the upper Hessian bound in \ref{as:hessian}.
\end{lemma}
\begin{proof}
From the definition of $m_{k}$,
\[
m_{k}(p^{\mathrm{C}}_{k})=f(x_{k})+g_{k}^{\top}p^{\mathrm{C}}_{k}
+\tfrac12(p^{\mathrm{C}}_{k})^{\top}B_{k}p^{\mathrm{C}}_{k}.
\]
Because $p^{\mathrm{C}}_{k}$ is colinear with $-g_{k}$,
\[
g_{k}^{\top}p^{\mathrm{C}}_{k}= -\|g_{k}\|\,\Delta_{k}.
\]
Using $\lambda_{\max}(B_{k})\le M$,
\[
(p^{\mathrm{C}}_{k})^{\top}B_{k}p^{\mathrm{C}}_{k}
\le M\|p^{\mathrm{C}}_{k}\|^{2}=M\Delta_{k}^{2}.
\]
Thus
\[
m_{k}(0)-m_{k}(p^{\mathrm{C}}_{k})
\ge \|g_{k}\|\Delta_{k} -\tfrac12 M\Delta_{k}^{2}
= \frac{\Delta_{k}}{2}\|g_{k}\| +\Bigl(\frac{\Delta_{k}}{2}\|g_{k}\|- \frac12 M\Delta_{k}^{2}\Bigr)
\ge \frac{\Delta_{k}}{2}\|g_{k}\| -\frac{M}{2}\Delta_{k}^{2}.
\]
\end{proof}

%%%%%%%%%%%%%%%%%%%%%%%%%%%%%%%%%%%%%%%%%%%%%%%%%%%%%%%%%%%%%%%%%%%%%%%%%%%%%%%%
\section*{Main Proof (Step‑by‑Step Derivation)}
We prove Theorem~\ref{thm:main}.  Throughout the proof we use the notation
$g_{k}:=\nabla f(x_{k})$, $B_{k}:=\nabla^{2}f(x_{k})$, and $\Delta_{k}$ as defined
above.

\medskip
\noindent\textbf{[Step 1]} \emph{Relating actual reduction to model decrease.}  
From the definition of $\rho_{k}$ (2) and Lemma~\ref{lem:model} (L1),
\[
\begin{aligned}
\rho^{\mathrm{act}}_{k}
&= f(x_{k})-f(x_{k}+p_{k})
   \\
&\ge m_{k}(0)-m_{k}(p_{k}) - \tfrac12 L\|p_{k}\|^{2}
   \qquad\text{(by Lemma~\ref{lem:model})}\\
&= \rho^{\mathrm{pred}}_{k} - \tfrac12 L\|p_{k}\|^{2}.
\end{aligned}
\tag{P1}
\]

\noindent\textbf{[Step 2]} \emph{Bounding the denominator of $\rho_{k}$.}  
Since $\rho_{k}= \rho^{\mathrm{act}}_{k}/\rho^{\mathrm{pred}}_{k}$,
using (P1) we obtain
\[
\rho_{k}
\ge 1 - \frac{L\|p_{k}\|^{2}}{2\,\rho^{\mathrm{pred}}_{k}} .
\tag{P2}
\]

\noindent\textbf{[Step 3]} \emph{Guaranteeing sufficient predicted decrease.}  
From the Cauchy condition (1) and Lemma~\ref{lem:cauchy} (L2),
\[
\begin{aligned}
\rho^{\mathrm{pred}}_{k}
&\ge \kappa_{\mathrm{c}}\bigl(m_{k}(0)-m_{k}(p^{\mathrm{C}}_{k})\bigr)\\
&\ge \kappa_{\mathrm{c}}\Bigl(\frac{\Delta_{k}}{2}\|g_{k}\|- \frac{M}{2}\Delta_{k}^{2}\Bigr)
\quad\text{(by Lemma~\ref{lem:cauchy})}.
\end{aligned}
\tag{P3}
\]

\noindent\textbf{[Step 4]} \emph{Choosing a lower bound on $\Delta_{k}$.}  
Define
\[
\Delta_{\min}:= \min\Bigl\{\Delta_{0},
\frac{(1-\eta_{1})\eta_{1}}{L}\Bigr\}.
\tag{P4}
\]
Standard trust‑region analysis (e.g., \cite{Conn2000}) shows that
as long as $\Delta_{k}\ge\Delta_{\min}$ the ratio $\rho_{k}$ cannot fall below
$\eta_{1}$.  Indeed, from (P2)–(P3) and the inequality $\|p_{k}\|\le\Delta_{k}$,
\[
\rho_{k}\ge
1-\frac{L\Delta_{k}^{2}}
      { \kappa_{\mathrm{c}} \bigl(\Delta_{k}\|g_{k}\| - M\Delta_{k}^{2}\bigr)}
\ge 1-\frac{L\Delta_{k}}{\kappa_{\mathrm{c}}\|g_{k}\| - M\Delta_{k}} .
\]
If $\Delta_{k}\le\frac{(1-\eta_{1})\eta_{1}}{L}$ the right‑hand side is at least
$\eta_{1}$.  Hence any iteration with $\Delta_{k}<\Delta_{\min}$ will trigger a
radius reduction (Step 5 below).  Consequently, after a finite number of
reductions the radius stabilises at $\Delta_{\min}$ and never falls below it.

\noindent\textbf{[Step 5]} \emph{Actual decrease when the step is accepted.}  
When $\rho_{k}\ge\eta_{1}$, the step is accepted and by definition
$\rho^{\mathrm{act}}_{k}= \rho_{k}\,\rho^{\mathrm{pred}}_{k}\ge\eta_{1}\rho^{\mathrm{pred}}_{k}$.
Using (P3),
\[
f(x_{k})-f(x_{k+1})
\;=\;\rho^{\mathrm{act}}_{k}
\;\ge\;\eta_{1}\kappa_{\mathrm{c}}
\Bigl(\frac{\Delta_{k}}{2}\|g_{k}\|- \frac{M}{2}\Delta_{k}^{2}\Bigr).
\tag{P5}
\]

\noindent\textbf{[Step 6]} \emph{Bounding the quadratic term.}  
Because $\Delta_{k}\le\Delta_{\max}$ and $\|g_{k}\|$ is non‑negative,
\[
\frac{M}{2}\Delta_{k}^{2}\le \frac{M}{2}\Delta_{\max}\Delta_{k}
\le \frac{1-\eta_{1}}{2}\,\Delta_{k}\|g_{k}\|
\quad\text{(choose }\Delta_{\max}\text{ small enough or use } \Delta_{k}\ge\Delta_{\min}).
\]
Thus from (P5),
\[
f(x_{k})-f(x_{k+1})
\;\ge\;
\frac{\eta_{1}\kappa_{\mathrm{c}}(1-\eta_{1})}{2}\,\Delta_{k}\|g_{k}\|.
\tag{P6}
\]

\noindent\textbf{[Step 7]} \emph{Linking radius to gradient norm.}  
From the radius update rule (4), whenever a step is accepted we have either
$\rho_{k}\ge\eta_{2}$ (radius enlarged) or $\eta_{1}\le\rho_{k}<\eta_{2}$
(radius unchanged).  In both cases $\Delta_{k+1}\ge\Delta_{k}$.  Hence after the
first iteration for which $\Delta_{k}\ge\Delta_{\min}$, the radius never
decreases again.  Consequently, for all subsequent accepted steps,
\[
\Delta_{k}\;\ge\;\Delta_{\min}.
\tag{P7}
\]

\noindent\textbf{[Step 8]} \emph{Summing the descent inequality.}  
Let $\mathcal{A}$ denote the set of indices where the step is accepted.
Summing (P6) over $k\in\mathcal{A}\cap\{0,\dots,K-1\}$ and using (P7) yields
\[
\sum_{k\in\mathcal{A}\cap\{0,\dots,K-1\}}
\|g_{k}\|
\;\le\;
\frac{2\bigl(f(x_{0})-f^{*}\bigr)}
     {\eta_{1}\kappa_{\mathrm{c}}(1-\eta_{1})\,\Delta_{\min}} .
\tag{P8}
\]

\noindent\textbf{[Step 9]} \emph{From sum of norms to minimum norm.}  
Since the left‑hand side of (P8) contains at most $K$ non‑negative terms,
\[
K\;\min_{0\le k<K}\|g_{k}\|
\;\le\;
\sum_{k\in\mathcal{A}\cap\{0,\dots,K-1\}}\|g_{k}\|
\;\le\;
\frac{2\bigl(f(x_{0})-f^{*}\bigr)}
     {\eta_{1}\kappa_{\mathrm{c}}(1-\eta_{1})\,\Delta_{\min}} .
\]
Rearranging gives
\[
\min_{0\le k<K}\|g_{k}\|
\;\le\;
\frac{2\bigl(f(x_{0})-f^{*}\bigr)}
     {\eta_{1}\kappa_{\mathrm{c}}(1-\eta_{1})\,\Delta_{\min}\,K}.
\tag{P9}
\]

\noindent\textbf{[Step 10]} \emph{Replacing $\Delta_{\min}$ by a Lipschitz constant.}  
From the definition (P4), $\Delta_{\min}\ge\frac{(1-\eta_{1})\eta_{1}}{L}$.
Plugging this bound into (P9) yields exactly the rate (5):
\[
\boxed{
\displaystyle
\min_{0\le k<K}\|\nabla f(x_{k})\|^{2}
\le
\frac{2L\bigl(f(x_{0})-f^{*}\bigr)}
     {\kappa_{\mathrm{c}}(1-\eta_{1})\eta_{1}\,K}.
}
\]

\noindent\textbf{[Step 11]} \emph{Limit argument.}  
Since the right‑hand side of (5) tends to $0$ as $K\to\infty$, the sequence
$\{\|\nabla f(x_{k})\|\}$ has a subsequence converging to $0$.
Because $f$ is convex, any accumulation point of $\{x_{k}\}$ is a global minimiser,
and the whole sequence converges to $x^{*}$ with the stated $O(k^{-1/2})$ rate.

This completes the proof of Theorem~\ref{thm:main}. \hfill $\square$

%%%%%%%%%%%%%%%%%%%%%%%%%%%%%%%%%%%%%%%%%%%%%%%%%%%%%%%%%%%%%%%%%%%%%%%%%%%%%%%%
\section*{Rate Derivation (With Constants)}
From (5) we can write an explicit iteration‑complexity bound:
to guarantee $\min_{0\le k<T}\|\nabla f(x_{k})\|\le\varepsilon$ it suffices that
\[
T\;\ge\;
\frac{2L\bigl(f(x_{0})-f^{*}\bigr)}
     {\kappa_{\mathrm{c}}(1-\eta_{1})\eta_{1}\,\varepsilon^{2}} .
\]
All constants on the right‑hand side are known a priori (or can be estimated),
so the algorithm enjoys a \emph{deterministic} worst‑case complexity of
$O(\varepsilon^{-2})$ gradient evaluations, matching the optimal rate for
first‑order methods on convex smooth problems.

%%%%%%%%%%%%%%%%%%%%%%%%%%%%%%%%%%%%%%%%%%%%%%%%%%%%%%%%%%%%%%%%%%%%%%%%%%%%%%%%
\section*{Special Cases}
\begin{enumerate}[label=\alph*)]
\item \textbf{Exact subproblem solution ($\kappa_{\mathrm{c}}=1$).}
The bound (5) simplifies to
\[
\min_{k<T}\|\nabla f(x_{k})\|^{2}
\le
\frac{2L\bigl(f(x_{0})-f^{*}\bigr)}
     {(1-\eta_{1})\eta_{1}\,T},
\]
which is the tightest possible under the given assumptions.

\item \textbf{Quadratic objective.}
If $f(x)=\tfrac12 x^{\top}Qx+ c^{\top}x$ with $Q\succeq 0$,
the model $m_{k}$ coincides with $f$ itself, i.e., $m_{k}=f$.
Thus $\rho_{k}=1$ for every iteration, the radius never shrinks, and the method
reduces to the classical Newton step with a safeguard $\|p_{k}\|\le\Delta_{k}$.
Convergence is achieved in a single iteration when $\Delta_{0}\ge\|Q^{-1}c\|$.

\item \textbf{Large‐scale problems (diagonal $B_{k}$).}
When $B_{k}$ is taken as a diagonal approximation of the Hessian,
the subproblem (TR$_k$) decouples into $n$ independent scalar trust‑region
problems, each solvable in closed form.  The same convergence guarantees hold
because Lemma~\ref{lem:model} only requires the global Lipschitz constant $L$,
not the exact Hessian.
\end{enumerate}

%%%%%%%%%%%%%%%%%%%%%%%%%%%%%%%%%%%%%%%%%%%%%%%%%%%%%%%%%%%%%%%%%%%%%%%%%%%%%%%%
\begin{thebibliography}{9}
\bibitem{Conn2000}
A.~R. Conn, N.~I.~M. Gould, and Ph.~L. Toint,
\emph{Trust‑Region Methods},
SIAM, 2000.
\end{thebibliography}

\end{document}